% ============================================================
%  III. WHOLE-BODY ACTIVE PERCEPTION AND CONTROL
% ============================================================
\section{Whole-Body Active Perception and Control}
\label{sec:active_perception}

The core contribution of \ac{TERESA} is a whole-body active perception
framework: the robot continuously adapts its posture---body height, pitch, and
yaw---to maximise perceptual quality. This section describes how each of the
four mission phases defined in Section~\ref{sec:problem_formulation} is
resolved: hybrid confidence-gated search (Phase~A), anticipatory approaching
with concurrent scanning (Phase~B), exposure body inspection with interactive
review (Phase~C), and \ac{FAST} ultrasound acquisition with compliant contact
(Phase~D).

The term \emph{whole-body} here denotes a perceptual extension of the classical
\ac{WBC} principle. Whereas traditional \ac{WBC} coordinates all degrees of
freedom to minimise kinematic error, \ac{TERESA} extends this coordination to
the perceptual domain: body posture---height, pitch, and yaw---is treated as an
explicit control resource that shapes the robot's sensory field. By commanding
posture to improve what the robot sees, at higher resolution and from more
informative viewpoints, \ac{TERESA} elevates the legged base from a passive
transport platform to an active participant in the perception-action loop. This
principle underpins every phase of the mission, from searching through
approaching to scanning.

\subsection{Phase A --- Searching}
\label{sec:phase_searching}

Before the robot can approach the patient, it must first localise them
reliably. In an unstructured emergency scenario, the patient's position is
unknown a priori. The \ac{EE}-mounted RealSense benefits from the Z1 arm's full
six-\ac{DOF} mobility, which can orient the camera in any direction without
moving the base. The body-mounted Orbbec, however, has a fixed orientation
relative to Spot: its downward field of view depends entirely on body pitch,
making pitch the only posture parameter that directly improves its detection of
a patient lying on the ground. \ac{TERESA} therefore avoids time-consuming base
rotations and instead searches by walking forward while sweeping body pitch.

\subsubsection{Forward search with pitch sweep}

The robot walks forward at low speed while the body pitch is adjusted through
a sequence of angles $\varphi \in [0, \varphi_{\max}]$. At each pitch, two
cameras observe the environment in complementary roles. The body-mounted Orbbec
runs YOLO11-pose and a posture classifier that outputs a confidence score
$c \in [0, 1]$. The \ac{EE}-mounted RealSense, actively oriented by the Z1
arm, runs a torso tracker that scans the surroundings independently of the
base heading. Two conditions are evaluated:

\begin{itemize}
  \item \textbf{Orbbec direct detection}: if the body-mounted camera detects a
        lying patient with high confidence at the current pitch, the robot
        transitions immediately to target locking.
  \item \textbf{RealSense guidance}: if the \ac{EE} camera detects a potential
        body while the Orbbec has not yet confirmed, the robot enters a
        semi-locking phase. It walks toward the detected torso and adjusts
        pitch to give the Orbbec a clean observation window.
\end{itemize}

If neither condition is met after the pitch sweep, the robot continues walking
forward and repeats the sweep, progressively covering the search area without
the idle time of discrete base rotations. This strategy exploits the arm's
native mobility for rotational search while using body pitch---the one viewing
parameter the arm cannot compensate for---to optimise the Orbbec's fixed
viewpoint.

\subsubsection{Dual-sensor confidence lock}

Patient localisation is confirmed through a hybrid lock that fuses information
from both cameras. When the Orbbec detects a lying posture with high
confidence, the robot enters a locking state: the arm returns to an elevated
home pose, and the coordinator collects a sequence of approach point estimates
in the world-fixed odometry frame $\mathcal{F}_W$, computing their mean as the
navigation target. Concurrently, the NLF pipeline produces a refined 3D
skeleton estimate by accumulating two valid detections with exponential moving
average smoothing, providing higher-fidelity anatomical data for the subsequent
approach phase. The transition to the approach phase is gated by both the
arm's home trajectory and the collection of the required number of samples. If
the Orbbec detection is temporarily lost, the system tolerates a brief absence
before resuming the search from the current position.

\subsection{Phase B --- Approaching}
\label{sec:phase_approaching}

In traditional mobile manipulation pipelines, the robot first navigates to a
target position, stops, and only then begins visual scanning---a sequential
approach that introduces significant idle time. \ac{TERESA} eliminates this
idle period by executing the body scan concurrently with the approach
navigation.

\subsubsection{Whole-body coordination}

During the approach, two independent controllers run in parallel. The Spot base
is driven by a proportional controller that computes $v_x$ and $\omega_z$ from
the distance and angle to the target in $\mathcal{F}_W$. The arm is controlled
by a \ac{WBC} formulation that orients the \ac{EE}-mounted RealSense toward the
patient's torso centroid. The joint velocities are computed via a damped
pseudo-inverse of the arm Jacobian:
\begin{equation}
\dot{\mathbf{q}} = \mathbf{J}_{\mathrm{arm}}^\top
\bigl(\mathbf{J}_{\mathrm{arm}} \mathbf{J}_{\mathrm{arm}}^\top
+ \mu\,\mathbf{I}_6\bigr)^{-1} \mathbf{v}_{\mathrm{des}},
\label{eq:damped_pinv}
\end{equation}
where $\mu > 0$ is a damping factor that prevents singularities near
kinematic limits, and $\mathbf{v}_{\mathrm{des}} \in \mathbb{R}^6$ is the
desired end-effector spatial velocity. The damped pseudo-inverse solves
$\mathbf{J}_{\mathrm{arm}} \dot{\mathbf{q}} = \mathbf{v}_{\mathrm{des}}$ in a
least-squares sense, mapping a commanded \ac{EE} motion to the joint space.
Here only the angular component of $\mathbf{v}_{\mathrm{des}}$ is active: it
encodes the rotation required to align the RealSense camera with the look-at
direction $\mathbf{x}_{EE} = (\mathbf{p}^* -
\mathbf{p}_{EE})/\|\mathbf{p}^* - \mathbf{p}_{EE}\|$, keeping the patient's
torso centred in the frame. The translational component of
$\mathbf{v}_{\mathrm{des}}$ is zero, as the arm holds a nominal home position
during approach. The base velocity is modulated by a
QualityMonitor that scales the approach speed in proportion to detection
confidence, avoiding premature commitment to uncertain targets while preventing
unnecessary stops.

\subsubsection{Adaptive anticipatory scan grid}

The arm executes a perceptual body scan using a dynamically generated Cartesian
grid of poses around the current \ac{EE} position. During the pre-approach
phase, the RealSense tracker publishes per-keypoint confidence scores for the
four torso landmarks (shoulders and hips), which drive the grid generation. If
all four keypoints are detected with high confidence, the scan uses a minimal
two-pose grid: the current \ac{EE} pose advanced toward the patient plus an
elevated viewpoint. If any keypoint falls below the threshold, the grid expands
to a three-pose comprehensive pattern with symmetric lateral offsets, an
elevated frontal viewpoint, and the home pose interleaved between lateral
transits. All scanning poses include a forward advance toward the patient. At
each pose, depth and skeleton data are fused into 3D keypoint estimates via a
weighted average accounting for detection rate and spatial stability. From the
fused torso model, the five canonical \ac{FAST} target poses are computed as
fixed offsets from the estimated torso centroid in $\mathcal{F}_0$.

\subsection{Phase C --- Exposure}
\label{sec:phase_exposure}

After the anticipatory scanner has localised the patient's anatomy, Phase~C
executes a systematic visual inspection of the body surface, serving the ``E''
(Exposure) component of the \ac{ABCDE} protocol. A key design principle is that
the operator remains in the loop: the system automates the scanning motion but
preserves the clinician's ability to inspect, re-examine, and override.

\subsubsection{Posture optimisation for body reconfiguration}
\label{sec:body_reconf}

Once the anticipatory body scanner has localised the patient's anatomy and
computed the target set $\mathcal{T}$ (both \ac{FAST} ultrasound targets and
exposure grid points) in $\mathcal{F}_W$, the system pre-plans an optimal body
posture for each target through a two-tier kinematically-aware stance search
with \ac{IK}-driven retry. This optimisation is shared with Phase~D
(Section~\ref{sec:phase_fast}).

The posture optimisation space spans three parameters: body height $h \in
[h_{\min}, h_{\max}]$ and pitch $\varphi \in [0, \varphi_{\max}]$, as defined
in Section~\ref{sec:problem_formulation}, plus a body-frame base displacement
$dy \in \mathbb{R}$ along the patient's head-to-foot axis, expressed in the
patient-aligned body frame (Section~\ref{sec:system_architecture}).
The discrete search grids are $\mathcal{H} = \{h_1, \dots, h_{N_h}\}$ and
$\mathcal{P} = \{\varphi_1, \dots, \varphi_{N_p}\}$, with $dy$ sampled
uniformly at \SI{0.10}{\metre} intervals.  The cost function includes a
regularisation term that penalises lateral displacement:
\begin{equation}
  J(h, \varphi, dy) =
  \bigl\|\mathbf{T}_{\mathcal{F}_W}^{\mathcal{F}_0}(h, \varphi, dy)\,
  \mathbf{p}^* - \mathbf{p}_{\mathrm{sweet}}\bigr\|
  + \lambda\,|dy|,
  \label{eq:posture_cost}
\end{equation}
where $\mathbf{p}_{\mathrm{sweet}} \in \mathbb{R}^3$ is the empirically
determined centre of the Z1 arm's preferred manipulation region in
$\mathcal{F}_0$, and $\lambda = 0.50$ weights lateral movement against
reachability.  The value $\lambda=0.50$ was chosen empirically to balance
reachability against energy cost: $\lambda>1.0$ prevented necessary base
displacement, while $\lambda<0.25$ caused excessive walking for marginal
reachability gains.  This penalty encourages the optimiser to prefer height and pitch
adjustments over base walking---reducing energy consumption and producing
smoother, more predictable motion---while still allowing longitudinal
displacement when strictly necessary.  The two-tier escalation proceeds as
follows:

\begin{enumerate}
  \item \textbf{2D search}---$(h, \varphi)$ only: a coarse grid sweeps
        $\mathcal{H} \times \mathcal{P}$ ($N_h N_p$ combinations, typically
        12).  For each candidate, the system transforms the target to
        $\mathcal{F}_0$ and queries the \ac{IK} solver.  If a posture yields
        a solution within a \SI{2}{\second} timeout, it is selected.
  \item \textbf{3D escalation}---$(dy, h, \varphi)$: when the 2D search fails,
        the optimiser adds longitudinal base displacement $dy$ (Spot walks
        along the patient's body), expanding to $N_{dy} N_h N_p$
        combinations.  If the \ac{IK} solver also times out at this tier, the
        point is skipped---the arm cannot reach it even with full base
        repositioning.
\end{enumerate}

The optimisation is computed once per target after $\mathcal{T}$ is available,
imposing no real-time computational burden.  Body reconfiguration is executed
between targets: Spot assumes the pre-computed posture $(h_i^*, \varphi_i^*)$,
settles, and the base navigator drives the $dy$ displacement before the Z1 arm
executes the \ac{IK} trajectory toward the transformed target.

\subsubsection{Posture-adaptive scan grid}

The scan grid adapts to the patient's estimated body pose. Four torso
keypoints---left and right shoulders, left and right hips---define a bilinear
patch over the body surface.  Let $\mathbf{p}_{SL}, \mathbf{p}_{SR},
\mathbf{p}_{HL}, \mathbf{p}_{HR} \in \mathbb{R}^3$ denote these keypoints in
$\mathcal{F}_W$.  The grid is generated by interpolating uniformly along the
two parametric directions: head-to-foot ($u$-axis, from shoulders to hips) and
lateral ($v$-axis, from left to right).  A grid point at parameters
$(u, v) \in [0,1] \times [0,1]$ is computed as
$\mathbf{p}(u,v) = (1{-}u)(1{-}v)\,\mathbf{p}_{SL} + u(1{-}v)\,\mathbf{p}_{HL}
+ (1{-}u)v\,\mathbf{p}_{SR} + uv\,\mathbf{p}_{HR}$, with a fixed height offset
to maintain a safe camera-to-skin distance.  Typical grid resolution is
\SI{0.15}{\metre} in both directions, producing 15--25 points depending on
torso dimensions.  When the NLF prior is available, all 24 \ac{SMPL} joints
are populated, extending the grid to the head, hands, and feet for full-body
coverage.  Grid points are visited in a serpentine path to minimise arm
travel.

\subsubsection{Injury detection}

At each grid point, Grounding DINO---a zero-shot
open-vocabulary object detector---runs on the RealSense RGB stream at
approximately 1\,fps on the Jetson Orin GPU, prompted with a vocabulary
of wound-related terms (wound, laceration, bleeding, burn, etc.). Detections
are projected to 3D world coordinates via the registered depth channel and
associated with their grid point index.

\subsubsection{Web-based interactive review}

Following the sweep, the system enters an interactive review phase: the web
interface overlays the exposure grid on the RealSense feed with detected
injuries highlighted as colour-coded bounding boxes.  The operator can click
any grid point to command the robot to return to that body region for
re-inspection.  The full re-positioning mechanism---posture optimisation,
settle, base navigation, and \ac{IK} trajectory replay---is detailed in
Section~\ref{sec:system_architecture}.  A termination command ends the review
and advances the mission to Phase~D.

\subsection{Phase D --- \acs{FAST}}
\label{sec:phase_fast}

Phase~D performs autonomous ultrasound acquisition at the $N$ \ac{FAST}
targets estimated during Phase~B. Body posture optimisation for each
$\mathbf{T}_i^* \in \mathcal{T}$ follows the same two-tier procedure described
in Section~\ref{sec:body_reconf}, applied to the \ac{FAST} target set. The
distinctive requirement of this phase is compliant probe contact with the
patient's body surface.

\subsubsection{Contact control}

The Unitree Z1 does not provide native impedance control, yet ultrasound
scanning requires compliant probe contact. \ac{TERESA} implements a custom
Cartesian impedance controller on the Z1 torque interface at \SI{500}{\hertz}.
The desired \ac{EE} pose is computed relative to the torso surface centroid
$\mathbf{p}_s$ and outward normal $\mathbf{n}_s$, placing the probe at
$\mathbf{p}_{\mathrm{des}} = \mathbf{p}_s + d\,\mathbf{n}_s$, where $d$ is the
desired contact offset.  Joint torques are computed as
\begin{equation}
\boldsymbol{\tau} = \mathbf{J}_{\mathrm{arm}}(\mathbf{q})^\top
\bigl[\mathbf{K}_p(\mathbf{p}_{\mathrm{des}} - \mathbf{p}_{EE})
- \mathbf{K}_d\,\dot{\mathbf{p}}_{EE}\bigr]
+ \mathbf{g}(\mathbf{q}) + \mathbf{C}(\mathbf{q}, \dot{\mathbf{q}}),
\label{eq:impedance}
\end{equation}
where $\mathbf{K}_p, \mathbf{K}_d \in \mathbb{R}^{3 \times 3}$ are adaptive
stiffness and damping gains that scale linearly with arm extension to prevent
overloading at the workspace boundary, $\mathbf{g}(\mathbf{q})$ and
$\mathbf{C}(\mathbf{q}, \dot{\mathbf{q}})$ are the gravity and Coriolis terms
computed via Pinocchio~\cite{carpentier2019pinocchio}, and
$\dot{\mathbf{p}}_{EE} = \mathbf{J}_{\mathrm{arm}}(\mathbf{q})\,
\dot{\mathbf{q}}$ is the end-effector velocity.  A safe switching service
alternates between the impedance controller and a position-controlled joint
trajectory controller, defaulting to position control on error.
